%% ============================================================
%%  Chapter 5 - Methodology and Evaluation Framework
%% ============================================================
\chapter{Methodology and Evaluation Framework}
\label{ch:methodology}

\section{Research Design}
\label{sec:design}

The study uses design-science systems research~\cite{hevner2004design}, controlled experimentation, emulation, simulation, and contextual validation. The artifact is the safety-constrained subscriber-centric broadband architecture. Evaluation tests whether the artifact satisfies authentication, isolation, QoS, accounting, edge feasibility, optimisation, LEO-assisted resilience, and Nepal-context deployment requirements.

Five evidence tracks are defined:
\begin{description}
  \item[Track 1 - Prototype:] OpenWrt + FreeRADIUS + VLAN/\texttt{tc} + WireGuard/VXLAN. Provides measured ground truth for authentication latency, slicing overhead, and AP resource consumption.
  \item[Track 2 - Edge overhead:] Hardware-in-the-loop benchmarking on constrained APs under slice saturation.
  \item[Track 3 - Emulation:] Mininet-WiFi topologies calibrated against prototype measurements.
  \item[Track 4 - Simulation and optimisation:] NS-3 and PyTorch Geometric scenarios evaluating graph-based macro policy.
  \item[Track 5 - Contextual validation:] Nepal scenario mapping, NTA regulatory compatibility, and deployment plausibility.
\end{description}

\section{Dataset Design}
\label{sec:dataset}

All records are classified as Measured, Emulated, Simulated, or Contextual (M/E/S/C). Claims cannot be extended beyond their evidence class.

\begin{longtable}{p{3.0cm}p{5.2cm}p{5.0cm}}
  \caption{Dataset matrices, variables, and analytical role.}\label{tab:dataset}\\
  \toprule
  \textbf{Matrix} & \textbf{Core Variables} & \textbf{Analytical Role} \\
  \midrule
  \endfirsthead
  \toprule
  \textbf{Matrix} & \textbf{Core Variables} & \textbf{Analytical Role} \\
  \midrule
  \endhead
  Identity-session matrix (M/E) & Subscriber pseudonym, realm, visited AP, home domain, EAP method, auth result, failure code, auth latency, session duration, accounting ID. & Tests whether identity operates as primary access-control primitive. \\
  Slice-performance matrix (M/E) & Slice ID, VLAN/VNI, rate cap, throughput, latency, jitter, loss, queue occupancy, tunnel overhead, SLA status. & Tests subscriber-specific QoS and slice enforcement. \\
  Edge-resource matrix (M) & CPU, RAM, IRQ rate, crypto throughput, encapsulation overhead, queueing cost, telemetry cost. & Determines AP feasibility under resource constraints. \\
  Radio-topology matrix (M/E/S) & AP ID, subscriber ID, RSSI, SNR, channel, retry rate, channel utilisation, edge weight, wall class. & Constructs graph states and path-loss validation data. \\
  Queueing matrix (M/E/S) & Arrival increments $A_k(t+\Delta)-A_k(t)$, service rate $C_k$, buffer threshold $B_k$, loss process increments $L_k^B$. & Estimates diffusion coefficient, overflow risk, and finite-buffer loss. \\
  Backhaul-resilience matrix (M/E) & Backhaul mode, RTT, jitter, packet loss, outage, LEO handover risk, terminal power, failover time. & Evaluates terrestrial and LEO resilience. \\
  Accounting-consistency matrix (M/E) & RADIUS counters, AP counters, gateway counters, byte mismatch, termination cause, offline signature. & Tests normal and degraded-mode auditability. \\
  Control-loop matrix (M/E/S) & GNN inference time, telemetry aggregation delay, policy update latency, epoch staleness, rollback rate. & Validates two-timescale separation. \\
  Contextual deployment matrix (C) & Location class, terrain class, backhaul type, power reliability proxy, regulatory constraint. & Maps Nepal-specific conditions to emerging-economy scenario classes. \\
  \bottomrule
\end{longtable}

\section{Data-Plane Tooling Hierarchy}
\label{sec:edge-tooling}

Commodity baseline enforcement uses Linux \texttt{tc}, \htb{}, \fqcodel{}, VLANs, and optionally VXLAN/WireGuard depending on the benchmark path. eBPF/XDP is not assumed on all APs; it is evaluated only on premium ARMv8/x86 nodes as an advanced benchmark track. This prevents the proposal from claiming low-end feasibility for mechanisms that may require stronger kernel, driver, and CPU support.

\section{System-Level Telemetry Throttle}
\label{sec:telemetry}

To prevent telemetry from competing with user-plane data, the system-level
telemetry interval is

\begin{equation}
\Delta t_{\mathrm{eff}}^{\mathrm{sys}}
=
\operatorname{clip}_{[\Delta t_{\mathrm{base}},\Delta t_{\mathrm{sys,max}}]}
\left(
\Delta t_{\mathrm{base}}
\exp
\left[
\kappa_B
\left[
1-
\frac{B_{\mathrm{curr}}}{B_{\mathrm{nominal}}+\epsilon}
\right]_+
+
\kappa_{\mathrm{CPU}}
\frac{CPU_j(t)}{CPU_j^{\max}}
\right]
\right),
\label{eq:system-telemetry-throttle-clean}
\end{equation}

where \([x]_+=\max(0,x)\). The coefficient $\kappa_B$ controls sensitivity to
backhaul degradation, while $\kappa_{\mathrm{CPU}}$ controls sensitivity to AP CPU
pressure. These are calibrated non-negative sensitivity coefficients and are not
required to sum to one. The positive-part operator prevents the telemetry controller
from increasing sampling frequency merely because transient throughput exceeds the
nominal baseline. This equation is separate from the LEO-specific backoff in
Eq.~\eqref{eq:leo-backoff-clean}.

\section{Experimental Baselines}
\label{sec:baselines}

\begin{enumerate}
  \item B1 - WPA-PSK router-centric baseline.
  \item B2 - Captive portal baseline.
  \item B3 - VLAN-only static policy baseline.
  \item B4 - Greedy AP/channel allocation baseline.
  \item B5 - DQN-style unconstrained learning baseline.
  \item B6 - Vanilla GNN without state-dependent safety filtering.
  \item B7 - Terrestrial-only backhaul and LEO-assisted backhaul variants reported separately.
\end{enumerate}

A configuration is accepted only if it improves the composite performance score
while satisfying all hard feasibility constraints:

\begin{equation}
Accept(\mathcal{M})=1
\iff
S(\mathcal{M})-S(\mathcal{B})>\Delta_{\min}
\land
\pi\in\Pi_{\mathrm{safe}}
\land
\mathbf{R}^{\mathrm{edge}}_j\le\mathbf{R}^{\max}_j.
\label{eq:acceptance-rule-clean}
\end{equation}

The composite score is

\begin{equation}
S(\mathcal{M})
=
\zeta_1T_{5p}(\mathcal{M})
-
\zeta_2L_{95p}(\mathcal{M})
-
\zeta_3V_{\mathrm{SLA}}(\mathcal{M})
+
\zeta_4J(\tilde{\mathbf{b}})
-
\zeta_5O_{\mathrm{edge}}(\mathcal{M}),
\label{eq:composite-score}
\end{equation}

where \(T_{5p}\) is fifth-percentile throughput, \(L_{95p}\) is
ninety-fifth-percentile latency, \(V_{\mathrm{SLA}}\) is the SLA violation
rate, \(J(\tilde{\mathbf{b}})\) is Jain fairness over normalised subscriber
entitlements, and \(O_{\mathrm{edge}}\) is measured edge-compute overhead.

This prevents the evaluation from accepting a configuration merely because it
improves average throughput while degrading tail latency, fairness, SLA
satisfaction, or commodity-edge feasibility.

\section{Ablation Study Plan}
\label{sec:ablation}

The full model is compared with ablated variants:
\begin{equation}
  \mathcal{M}_{full},\;\mathcal{M}_{-graph},\;\mathcal{M}_{-policy},\;\mathcal{M}_{-physics},\;\mathcal{M}_{-adaptive},\;\mathcal{M}_{-surv},\;\mathcal{M}_{-LEO}.
  \label{eq:ablation}
\end{equation}

\section{Evaluation Metrics}
\label{sec:metrics}

\begin{table}[H]
\centering
\caption{Metric-to-research-question alignment.}
\label{tab:metrics}
\begin{tabularx}{\textwidth}{llX}
\toprule
\textbf{RQ} & \textbf{Metric Class} & \textbf{Operational Definition} \\
\midrule
RQ1 & Identity and policy & Authentication success, EAP/RADIUS latency, realm resolution success, policy-attribute correctness, reject accuracy. \\
RQ2 & Slice enforcement & Cross-slice leakage, tunnel overhead, per-user bandwidth compliance, shaped latency, accounting mismatch. \\
RQ2 & Edge feasibility & CPU, RAM, IRQ, crypto, tunnel, queueing, telemetry overhead. \\
RQ3 & Optimisation & Aggregate throughput, fifth-percentile throughput, Jain fairness, SLA violation, interference-cost reduction. \\
RQ3 & Queue stability & Estimated $\hat{\sigma}_{k,	au}^2$, finite-buffer loss rate, auxiliary $P(Q_k^\infty>B_k)$ estimate. \\
RQ1--3 & Deployment viability & Configuration complexity, AP density, terrestrial/LEO sensitivity, NTA compatibility. \\
\bottomrule
\end{tabularx}
\end{table}

\section{Statistical Analysis Plan}
\label{sec:stats}

For paired comparisons, the treatment effect is
\begin{equation}
  \Delta Y=\frac{1}{R}\sum_{r=1}^{R}(Y_{r,c^*}-Y_{r,c_0}).
  \label{eq:treatment-effect}
\end{equation}
A mixed-effects model accounts for AP, subscriber, and scenario variation:
\begin{equation}
  Y_{ijst}=\beta_0+\beta_1Config_i+\beta_2Density_j+\beta_3Backhaul_k+\beta_4Mobility_l+u_s+\varepsilon_{ijst}.
  \label{eq:mixed-effects}
\end{equation}

Because the evaluation compares multiple baselines, scenario classes, and metrics,
confirmatory hypothesis tests are corrected using the Holm--Bonferroni family-wise
error-rate procedure within each metric family. Exploratory secondary analyses may
additionally report Benjamini--Hochberg false-discovery-rate adjusted values, clearly
labelled as exploratory. The correction method is fixed before comparative baseline
results are analysed.
